\section{Further detail}
\label{app:detail}

The material below was moved out of the body for length. Nothing in it is altered.

\subsection{The household credit engine}

\emph{Supporting Section~\ref{sec:data}, on how a displacement level becomes household credit losses.}

The engine maps a displacement level onto household credit losses without naming a set of
displaced households. Displacement is a probability applied to every household in the working
core, defined as a household with at least one employed member aged 25 to 64. The share of
earners displaced is set equal to the share of the wage bill displaced, and each household's
earners are drawn against that share.

A household that loses an earner does not default; it becomes more likely to default. The
uplift is five percentage points for one displaced earner and eight for two
\citep{GerardiHerkenhoffOhanianWillen2018}, and losses are exposure at default times loss given
default, with severities of 0.25 to 0.40 on first-lien mortgages, 0.80 to 1.00 on cards, 0.45 to
0.65 on auto and 0.75 to 1.00 on student debt. Applying a portfolio loss rate to exposure at
default instead would count the probability of default twice.

Survey balances are benchmarked to official aggregates first, and two ratios do different work.
The \emph{under-reporting factor} is the official aggregate over the full survey universe and is a
survey correction; the \emph{coverage share} is the working-core balance over that universe and is
a fact about which households the analysis covers. Table~\ref{tab:underreporting} gives both. The
identity between them holds to three decimals on every book, and all four factors were rebuilt
independently.

One scope statement travels with every credit number here. The engine is first round only: house
prices, consumer spending, business revenue and the employment of everyone not displaced are held
fixed. Against a supervisory scenario, which moves all of those together, the figures are a lower
bound rather than an estimate.

\subsection{The fiscal module and the second round}

\emph{Supporting Section~\ref{sec:data}, on the two modules that carry displacement into the budget and into bank losses.}

The fiscal module converts displacement into a dollar amount against the wage bill and applies a
loss per dollar with three parts: the labor tax no longer collected on wages that are not
re-earned, less the capital tax collected on the income that displaced them, plus any outlay
response. The share of wages re-earned is governed by a reemployment rate fitted on the
displacement vintages against labor market slack and solved as a fixed point, since displacement
changes the slack the fit is a function of. Losses are reported for the terminal year rather than
as a cumulative total divided by a horizon, and they are not linear in displacement: the share
re-earned falls as displacement rises. Trust fund effects use each fund's own payroll share of its
own income.

The second round is the weakest module evidentially and is labeled scenario wherever it appears.
The net fall in wage income becomes a demand shortfall through the difference between the
propensity to consume out of labor and out of capital income
\citep{MianStraubSufi2020,FagerengHolmNatvik2021}; the shortfall is expressed as a share of output
and compared with the fall in output in the supervisory severely adverse scenario; and the
supervisor's own published loss totals are scaled by that ratio. House prices enter through an
income elasticity \citep{HarterDreiman2004}, and employment through an Okun coefficient
\citep{BallLeighLoungani2017}. Three elements have no source: the shape of the loss curve beyond
the supervisory point, for which we carry a linear, a capped and a convex mapping; the propensity
ranges, measured on transitory shocks where displacement is persistent; and the path itself. We
report the whole grid of 648 combinations as a band and never a point inside it.

\subsection{Institution-level data}

\emph{Supporting Section~\ref{sec:data}, on the balance-sheet panel and the capital tests.}

The second round is applied across individual balance sheets rather than to an aggregate. The
panel is every insured bank and credit union at 30 June 2026, with the loan book by category and
the capital measure each is tested against: the tier 1 leverage ratio against its four percent
threshold for banks, the net worth ratio against its six percent threshold for credit unions.
After exclusions it is \result{NInstitutions} institutions, \result{NBanks} banks and
\result{NCreditUnions} credit unions, holding \result{SystemAssetsBn} billion dollars of assets.

\result{ExcludedInstitutions} institutions holding \result{ExcludedAssetsBn} billion dollars are
excluded because they report neither tier 1 capital nor total equity, being overwhelmingly insured
US branches of foreign chartered institutions whose capital sits at the parent; a capital test
shows them breaching before any loss, which is an artifact of a missing field. A capital test is
only informative if almost nothing fails it before any loss, and after the exclusion
\result{BaselineBreaches} institutions breach at baseline, \result{BaselinePctAssets} percent of
system assets. Losses are allocated pro rata by category holdings, so dispersion reflects business
mix alone and is a lower bound on true dispersion.

\subsection{The two conventions in the accounts}

\emph{Supporting Section~\ref{sec:accounts}, on the two conventions the accounts rest on.}

\paragraph{Government debt enters by its receipts, not by its obligor.} Treasury debt is not
wage-backed because a household owes it, but to the extent that wage taxes service it, so it takes
the labor-linked share of federal receipts, social insurance contributions in full plus the wage
share of the individual income tax, measured at \result{ReceiptsPctLow} percent; state and local
debt is treated the same way on its own receipts mix. Without this convention a government funded
entirely from payroll taxes and one funded entirely from capital taxes would score identically.

\paragraph{The one-step rule is a convention, and it is the largest judgment in the paper.} We
trace servicing one step and stop, so a claim serviced out of business revenue takes zero direct
backing even though some of that revenue began as wages. Multifamily mortgages are the one
property class that does not take the zero, because the rent that services them is paid out of
wages with no firm in between. The rule is the first-round boundary applied to the liability side.
Tracing further is what the indirect variant does, and we report both: on the whole claim stock,
moving business classes to their measured indirect share raises the ratio by
\result{OneStepMoveAllClaimsPct} percent, more than every other judgment call in the accounts
together.

\subsection{The series, 1952 to 2025}

\emph{Supporting Section~\ref{sec:accounts}, on the path of the accounts since 1952.}

Figure~\ref{fig:series} runs the accounts back to 1952, and two things in it move differently.

The ratio itself is close to flat across seventy-three years, between
\result{DebtOnlySeriesLow} and \result{DebtOnlySeriesHigh} percent on the debt-only measure with
no trend. The share of the national debt stock that wages service is not a recent development.

The federal share of it traces a U: \result{SovUnion1952} percent in 1952,
\result{SovUnion1970} percent in 1970, \result{SovUnion2025} percent in 2025. The 1952 level is
almost entirely wartime Treasury debt, the debtor position alone standing at
\result{SovObligor1952} percent. It falls to 1970 as that debt shrinks against a growing claim
stock, then climbs in two datable steps: the agency book grows between 1970 and 1990, taking the
held-or-guaranteed position from \result{SovHeld1970} to \result{SovHeld1990} percent, and federal
debt grows between 2008 and 2020, taking the debtor position from \result{SovObligor2008} to
\result{SovObligor2020} percent.

Said plainly, because it qualifies the headline: a large part of the recent rise is growth in
federal debt itself rather than new exposure to households. The held-or-guaranteed position peaked
at \result{SovHeld2020} percent in 2020 and has since fallen to \result{SovHeld2025} percent,
while the whole net increase since 2008 sits in the debtor position.

\subsection{Households, by position in the wage distribution}

\emph{Supporting Section~\ref{sec:accounts}, on the household part of the wage-backed stock.}

Splitting the household part of the stock by wage quintile gives the gradient: wage-backed claims
per dollar of wages fall as pay rises, from \result{QTwoPerDollar} in the second quintile to
\result{QTopPerDollar} at the top, while the top quintile earns \result{QTopWageShare} percent of
the wage bill and carries \result{QTopClaimShare} percent of the household claims. Claims are
spread far more evenly than the income that services them, which is the sense in which a wage
shock is regressive before any policy response is chosen. Appendix Table~\ref{tab:quintile} gives
the detail.

The magnitude for the bottom quintile is withdrawn: an independent rebuild landed a factor of 2.5
to 3.4 away, and we cannot show its reading was wrong, because our own specification defined
neither the universe nor the normalization of that cell. The direction was reproduced and stands.

\subsection{Checks, and what this module does not reach}

\emph{Supporting Section~\ref{sec:fiscal}, on the published figures the assembled rate is checked against and the limitation that runs against our own finding.}

Our shareholder layer at central values, \result{ShareholderLayerPct} percent, corroborates the
Congressional Research Service's text figure of about \result{CRSTextFigure} percent for the
effective capital gains rate on corporate profits. It does not reproduce that report's published
band of \result{CRSBandLow} to \result{CRSBandHigh} percent as printed, and should not: the band
already carries a taxable-share adjustment over domestically held stock. Rescaled to a common base
it is \result{CRSRescaledLow} to \result{CRSRescaledHigh} percent, and our figure sits inside it.
We report the correction rather than the original comparison, because as first written our own
check failed against our own construction.

Foreign holders are treated as bearing no tax at the owner level, though dividends to them bear
withholding tax at treaty rates. The effect on the assembled rate is upward and we have not
quantified it; since our result is that the rate falls short, that is a limitation running against
our own finding. The debt share is an economy-wide stock ratio, \result{DebtShareLow} to
\result{DebtShareHigh} percent, standing in for a marginal flow; the evidence in
Section~\ref{sec:bust} points below that range, and because the debt-financed term is negative, a
lower debt share would raise the assembled rate.

\subsection{The federal share, by displacement level and by classification}

\emph{Supporting Section~\ref{sec:displacement}, on how the federal share moves with the displacement level.}

The federal share of the first round is not one number, and reporting it as one hides that it
rises with the displacement level. It also depends on a classification choice about the housing
agencies, and both readings travel together. Under the \emph{narrow} reading only the agency loss
beyond the agencies' own capital and a year of earnings counts as federal. Under the
\emph{conservatorship} reading the retained agency loss counts as federal, because the agencies
are in conservatorship and no private shareholder is positioned to absorb anything; only what
private credit cover actually pays is private.

At the ten percent level the federal share is \result{FedShareTenNarrowLow} to
\result{FedShareTenNarrowHigh} percent under the narrow reading and \result{FedShareTenConsLow} to
\result{FedShareTenConsHigh} percent under the conservatorship reading.
Figure~\ref{fig:fedshare} gives the whole path: it rises with displacement to a peak near the
fifty percent level and falls back, because the fiscal component saturates while credit losses
keep growing. The rise between the five and twenty-five percent levels was reproduced
independently, the rebuild's own result sitting inside our range at both ends.

\subsection{The housing agencies absorb their share without reaching the Treasury}

\emph{Supporting Section~\ref{sec:displacement}, on the agency book rebuilt from the Enterprises' filings.}

The agency book was rebuilt loan class by loan class from the Enterprises' own filings
\citep{FannieMae10K2025,FreddieMac10K2025}, which is what makes the classification question
answerable rather than assumed. Three measured results follow. A large part of each single-family
book carries no credit enhancement at all, \result{UnenhancedLow} to \result{UnenhancedHigh}
percent across the two Enterprises. Private cover transfers only \result{PrivateCoverLow} to
\result{PrivateCoverHigh} percent of an agency loss at the ten percent level, because the
structures are written for idiosyncratic losses and an event of this size passes most of the loss
to the guarantor. And the Treasury draw is \result{TreasuryDrawBn} at every level we ran:
\result{AgencyCapitalBn} billion dollars of capital stands against a maximum retained loss of
\result{MaxRetainedAgencyLossBn} billion. That is a first-round result with house prices fixed,
and a house price fall would change the severity on every loan in the book.

\subsection{Trust funds}

\emph{Supporting Section~\ref{sec:displacement}, on the payroll basis of the trust fund calculation.}

Trust fund effects are computed on a payroll-only basis: each fund's loss is the displaced share
of its own covered payroll, adjusted for wages that return, times that fund's own payroll share of
its own income. The two funds have different shares and different caps, the social security funds
being capped at the contribution base while the hospital insurance fund is not, which is why the
two rows in Table~\ref{tab:incidence} move differently across exposure groups. One consequence
carries into Section~\ref{sec:institutions}: replacement income funded from a capital tax cannot
repair these funds, because it is not covered wages.

\subsection{The exposures this engine does not reach}

\emph{Supporting Section~\ref{sec:banks}, on exposures the balance-sheet panel does not cover.}

Four sets of institutions are exposed in ways the panel does not capture, and naming them is more
useful than putting a number on them the data cannot support. \emph{State and local government} is
exposed through its own tax base, \result{StateLocalWageLinkedBn} billion dollars of wage-linked
income tax, \result{StateLocalWageLinkedShare} percent of its own tax receipts, and cannot run the
deficit through because of balanced-budget requirements; its in-year spending cut feeds the same
demand channel that produces nine tenths of the bank losses above, and that loop is not in the
engine, so its absence makes our figures conservative in a direction we can sign but not size.
\emph{Pension funds} are exposed on both sides at once, benefits fixed in nominal terms while
contributions are a share of covered payroll, so displacement widens the funding gap mechanically;
the accounts capture only the \result{PensionHoldingsBn} billion dollars they hold.
\emph{Insurers and private credit} hold both sides, insurers \result{InsurerHoldingsBn} billion
dollars of wage-backed claims against long-dated liabilities, while private credit cannot be
separated from the other-financial aggregate, so the \result{OtherFinancialHoldingsBn} billion
attributed to that sector is an upper bound by an unknown margin. And everything outside the
insured depository system, nonbank servicers and consumer lenders among them, is absent from the
panel; as a class they are more thinly capitalized than the institutions we measure, so the system
figures are again a lower bound.

\subsection{What is needed at each regime}

\emph{Supporting Section~\ref{sec:institutions}, on the instrument set each regime calls for.}

For each regime we state the minimum instrument set identified by the measured mechanisms, which
is a narrower claim than a policy recommendation: it is what the mechanisms we can measure imply,
and nothing more.

\paragraph{Inside the data, about ten percent of the wage bill.} The capital tax instrument, the
trust fund contribution base, the student loan trigger, and better measurement. The housing
agencies absorb this level without a Treasury draw and the banking system barely moves, so
mortgage instruments and ownership stakes are not needed here and a supervisory scenario is a
disclosure requirement rather than a binding constraint. This is the only regime fully inside the
observed data, and the one where we can say what is unnecessary as confidently as what is not.

\paragraph{Larger displacement, twenty-five to fifty percent.} Nearly everything, with the
composition changing rather than the scale. The break-even capital tax rate, \result{BreakEvenALow}
to \result{BreakEvenAHigh} percent with output preserved and \result{BreakEvenBLow} to
\result{BreakEvenBHigh} percent with output falling, reaches the top of the range anyone has
measured, so ownership and direct-claim instruments enter here precisely because the tax
instrument is at its limit. All of it is scenario and extrapolates the reemployment
relationship.

\paragraph{Sudden displacement.} The same set reordered: only contract instruments act inside a
quarter, so pre-authorization is the whole instrument. Tax changes and ownership stakes take
legislative time the regime does not have, and pre-funding is useless once the shock has arrived.

\paragraph{Near-total displacement.} The share of wages re-earned falls toward zero, so the
required capital tax rate approaches the whole effective labor tax rate: \result{BreakEvenALow} to
\result{BreakEvenAHigh} percent with output preserved and \result{BreakEvenBLow} to
\result{BreakEvenBHigh} percent with output falling, and higher still once the pass-through
coefficient of Section~\ref{sec:fiscal} is applied. Those rates lie far outside anything observed
in a modern tax system, which is a statement about precedent rather than about arithmetic:
capital taxation is not excluded in principle at those levels, but nothing in the record tells us
what a rate of that size would do to the base it is levied on. Ownership and direct-claim
instruments therefore become the more plausible channel, not the only possible one. This is where
the question stops being financial regulation and becomes a question about ownership.

\paragraph{The failure case.} Instruments that protect the wage tax base do nothing here.
Concentration limits on AI-linked lending and the capital tax rate are what remain, and the second
is the same row that appears in every displacement regime.

\subsection{What each type of institution should change}

\emph{Supporting Section~\ref{sec:institutions}, on what each institution type should change and what the evidence does not support.}

Each item is tied to a measured exposure in Section~\ref{sec:banks}, and each carries what the
evidence does not support, because the second half keeps the first half honest. Card-heavy lenders
should hold capital against the card book for a displacement scenario specifically, and this is
not a tail risk: it is the first model to break and it breaks inside the range we model. Credit
unions are undiversified by design rather than small and therefore safe. Large and regional banks
need a supervisory scenario routed through the second round, not tighter household underwriting,
which does not touch their exposure. Auto lenders should write displacement-contingent payment
clauses into contracts and securitization documentation, though not at moderate displacement,
where they show no breaches. Housing agencies should pre-authorize forbearance and raise credit
enhancement on the unenhanced book, not recapitalize, since no Treasury draw appears at any level
we ran. Insurers should report AI-linked and wage-linked holdings together; forced-sale rules are
not supported, since they hold to maturity against matched liabilities. For private credit the
instrument is disclosure sufficient to separate it from the aggregate it sits inside. Pension
funds should report the funding gap against a displacement path, and state and local governments
should build reserves against one.

\subsection{What the evidence does not support}

\emph{Supporting Section~\ref{sec:institutions}, on two claims that do not survive the measurement.}

Two claims are tempting and neither survives the measurement. \emph{Wholesale change to household
mortgage underwriting at moderate displacement} fails on three grounds: the flat debt service
result is partly mechanical, since underwriting caps debt service to income; physically exposed
households are less indebted on most measures we can compute; and mortgage portfolio lenders are
among the least exposed institution types.

\emph{Occupation-based underwriting} fails on a legal ground as well as an empirical one, and the
statement has to be careful because a looser version of it was wrong. Occupation is not a
prohibited basis under the Equal Credit Opportunity Act, whose list of prohibited bases does not
include it, and a creditor may consider information not used to discriminate on a prohibited
basis. The exposure is different in kind: because occupation correlates with protected classes,
occupation-based underwriting in housing credit exposes a lender to a discriminatory-effects claim
under the Fair Housing Act, on which the lender must carry a legally sufficient justification
under a burden-shifting framework. Combined with the finding that the exposure is spread almost
evenly across borrowers, so screening on it buys a lender little, the risk-adjusted case for
occupation-based underwriting is weak and the case for portfolio-level disclosure is strong.

\subsection{Emerging markets}

\emph{Supporting Section~\ref{sec:institutions}, on why the answer differs outside a reserve-currency economy.}

Three facts change the answer outside a reserve-currency economy and they compound. Imported AI
capital accrues its taxable surplus abroad, so the domestic rate is lower still and the tax
instrument is weakest where it is needed most. A sovereign that cannot borrow in a currency it
issues cannot absorb the shock the way the United States can, leaving pre-funding, which must be
built beforehand. And many such economies are already on a steep path: our twenty-year baseline
reaches \result{EMBaseline} percent of output before any displacement, with an increment at the
ten percent level of \result{EMIncrementLow} to \result{EMIncrementHigh} points against
\result{ReserveIncrementLow} to \result{ReserveIncrementHigh} for a reserve-currency issuer, a
ratio of about \result{EMRatio}. The measurement itself is unavailable there: for the one case we
attempted, the flow of funds data have long lags and incomplete household coverage, there is no
usable counterparty matrix, and high informality places much labor income outside both the tax
base and the formal credit system, which is the variable the statistic turns on.

\subsection{State-contingent triggers}

\emph{Supporting Section~\ref{sec:institutions}, on the triggers and their current values.}

Instruments are more useful attached to a trigger carrying its current value and a label saying
whether the threshold sits inside the range of the data behind it. The capital tax trigger is already tripped:
the assembled rate is \result{TauKBarkai} percent against a threshold of \result{RequiredEasy}
percent, both sides observed today. Prime-age nonemployment stands at \result{PrimeAgeNonemp}
against the observed maximum of \result{PrimeAgeMax}, where the reemployment fit leaves the data,
so crossing it is the moment the rest of this architecture can no longer be sized. The
debt-financed share of AI capital spending is \result{DebtFinancedShare} percent against the
twenty percent read off the equity-led episode, and is a lower bound. The combined social insurance
depletion date is \result{OASDIDepletion}, which arrives before any plausible displacement path,
so that instrument binds soonest for reasons unrelated to AI. Recent graduate unemployment at
\result{GradUnemployment} percent and underemployment at \result{GradUnderemployment} percent sit
in somebody else's data and speak to the channel our engine cannot see.

One trigger matters more than the others and is a statement about knowledge rather than the
economy. Past a displacement level of about forty-six percent of the physically exposed wage bill,
the reemployment construction has no admissible solution and no reemployment-based instrument can
be sized from these data. The architecture does not become wrong there; it becomes unsizeable.
Appendix~\ref{app:instruments} gives the full instrument table.

\subsection{Avenues for further work}

\emph{Supporting Section~\ref{sec:limitations}, one avenue for each limitation.}

Each addresses a limitation above. Build the
off-balance-sheet side from securitization prospectuses, vehicle-level filings and supervisory data
on lending secured against computing hardware, so the AI side can be measured rather than bounded
(L1). Separate realized gains from wage income inside the individual income tax base and re-run
the fiscal magnitudes (L2). Run a systematic search with a published protocol for prior
constructions of a labor backing decomposition (L3). Identify the holder remainder from loan-level
sources (L4). Replace the single labor tax rate with rates by position in the income distribution
(L5). Construct both bases for the taxable-holder share from the same holder data and report the
rate on each (L14). Extend the reemployment relationship using a source that observes larger
displacement episodes, which would turn the bands above a quarter of the wage bill into estimates;
let house prices, spending and non-displaced employment move together, which would also let the
state and local feedback loop be modeled rather than named; run a rebuild with an enforced blind
and expected values held by a third party; and quantify withholding tax on dividends to foreign
holders, reporting the corrected rate even though it runs against our own finding.

